% !TEX program = pdflatex
% ------------------------------------------------------------
% MATH 431 Homework Template (plug-and-chug friendly)
%
% Goals:
% - Same clean header style as your MATH 521 template.
% - Easy “Problem -> Solution” workflow.
% - Lightweight (no tcolorbox); uses mdframed for optional boxes.
% - Good for probability homework: sample space, probability measure, events.
%
% Quick use:
%   1) Edit the Metadata block.
%   2) For each exercise, use \problem{<section>.<number>} and paste the prompt.
%   3) Write your solution in the solution environment.
% ------------------------------------------------------------

\documentclass[12pt]{article}

% =========================
% 0) Metadata (EDIT ME)
% =========================
\newcommand{\CourseCode}{MATH 431}
\newcommand{\CourseTitle}{Introduction to Probability}
\newcommand{\CourseSection}{LEC 002} % change to 001 or 003 if needed
\newcommand{\Semester}{Spring 2026}
\newcommand{\Instructor}{Sarah Tammen} % LEC 001: Ander Aguirre | LEC 002: Sarah Tammen | LEC 003: Nicholas Derr
\newcommand{\MeetingInfo}{MWF 9:55--10:45, 6102 Sewell Social Sciences} % LEC 001 default; update if in 002/003
\newcommand{\HomeworkNumber}{6}
\newcommand{\YourName}{Alejandro De La Torre}
\newcommand{\DueDate}{March 6, 2026} % or e.g. February 2, 2026

% =========================
% 1) Core math tools
% =========================
\usepackage{amsmath, amssymb, amsthm}

% =========================
% 2) Lists and spacing
% =========================
\usepackage{enumitem}
\setlist[enumerate,1]{label=\textbf{(\alph*)}, leftmargin=2em, itemsep=0.25em}
\setlist[itemize]{leftmargin=1.5em, itemsep=0.25em}
\setlength{\parskip}{0.35em}
\setlength{\parindent}{0pt}

% =========================
% 3) Page geometry + header/footer
% =========================
\usepackage{geometry}
\geometry{margin=1in}

\usepackage{fancyhdr}
\pagestyle{fancy}
\fancyhf{}
\lhead{\CourseCode\ --- \CourseSection, \Semester}
\rhead{Homework \HomeworkNumber}
\cfoot{\thepage}
\renewcommand{\headrulewidth}{0.4pt}
\setlength{\headheight}{15pt}
\addtolength{\topmargin}{-2pt}

% =========================
% 4) Links (subtle)
% =========================
\usepackage[hidelinks]{hyperref}
\setcounter{tocdepth}{1}

% =========================
% 5) Figures (optional)
% =========================
\usepackage{graphicx}
\graphicspath{{images/}}

% =========================
% 6) Lightweight boxes (optional)
% =========================
\usepackage{mdframed}
\newmdenv[
  skipabove=0.6\baselineskip,
  skipbelow=0.6\baselineskip,
  linewidth=0.6pt,
  roundcorner=4pt,
  innerleftmargin=10pt,
  innerrightmargin=10pt,
  innertopmargin=8pt,
  innerbottommargin=8pt
]{boxnote}

% =========================
% 7) Probability / notation macros
% =========================
\newcommand{\R}{\mathbb{R}}
\newcommand{\N}{\mathbb{N}}
\newcommand{\Z}{\mathbb{Z}}

\newcommand{\OmegaSpace}{\Omega}
\newcommand{\FField}{\mathcal{F}}
% Probability measure in case it is relevant or want to distinguish it from P(A)
%\newcommand{\PMeasure}{\mathbb{P}} 
\newcommand{\E}{\mathbb{E}}

\newcommand{\given}{\,\middle|\,}
\newcommand{\ind}{\mathbf{1}}

% Common “defined as” symbol
\newcommand{\defeq}{\vcentcolon=}

% =========================
% 8) Problem command (TOC + PDF bookmarks)
% Usage:
%   \problem{<label>}
%   \problem[Short title]{<label>}
% Example:
%   \problem{1.4}
%   \problem[Sample space]{1.4}
% =========================
\makeatletter
\newcommand{\problem}[2][]{%
  \def\prob@title{Problem #2\if\relax\detokenize{#1}\relax\else\ (\textit{#1})\fi}%
  \phantomsection%
  \pdfbookmark[1]{\prob@title}{prob#2}%
  \addcontentsline{toc}{section}{\prob@title}%
  \section*{\prob@title}%
  \vspace{-0.25em}\hrule\vspace{0.8em}%
}
\makeatother

% =========================
% 9) Solution environment
% =========================
\newenvironment{solution}{%
  \noindent\textbf{Solution.}\ \ignorespaces
}{\par}

% Optional scratch / planning block
\newenvironment{scratch}{%
  \begin{boxnote}\textbf{Scratch / Setup.}\ \small
}{\end{boxnote}}

% Final answer command: expects math only (no $$ or \[ \])
\newcommand{\finalanswer}[1]{%
  \par\medskip
  \noindent\textbf{Final:}\ \fbox{$#1$}\par
} % AGAIN: MATH ONLY INSIDE THE BOX; if word answer, use \text{} inside or use \fbox{TEXT} or \framebox

% =========================
% 10) Title block
% =========================
\title{Homework \HomeworkNumber}
\author{\YourName \\
\CourseCode\ --- \CourseTitle \\
\CourseSection \\
Instructor: \Instructor}
\date{Due: \DueDate}

\begin{document}
\maketitle

% \section*{Collaboration / Resources}
% (Optional) Write “I worked alone” or list collaborators/resources.

\tableofcontents
\bigskip
\hrule
\bigskip
\newpage

% ============================================================
% Problems (duplicate blocks as needed)
% ============================================================


\problem{3.16}
Let $Z$ have the following density function
\[
f_Z(z)=
\begin{cases}
\frac{1}{7}, & 1 \le z \le 2,\\
\frac{3}{7}, & 5 \le z \le 7,\\
0, & \text{otherwise}.
\end{cases}
\]
Compute both $E[Z]$ and $\mathrm{Var}(Z)$.

\begin{solution}
\begin{scratch}
Support: $z\in[1,2]\cup[5,7]$, with $f_Z(z)=\frac{1}{7}$ on $[1,2]$ and $f_Z(z)=\frac{3}{7}$ on $[5,7]$.
Plan: compute $E[Z]=\int z f_Z(z)\,dz$ and $E[Z^2]=\int z^2 f_Z(z)\,dz$ over those intervals, then $\mathrm{Var}(Z)=E[Z^2]-(E[Z])^2$.
\end{scratch}

\[
E[Z]=\int_{-\infty}^{\infty} z f_Z(z)\,dz
= \frac{1}{7}\int_{1}^{2} z\,dz + \frac{3}{7}\int_{5}^{7} z\,dz
= \frac{1}{7}\cdot\frac{2^2-1^2}{2} + \frac{3}{7}\cdot\frac{7^2-5^2}{2}
= \frac{75}{14}.
\]

\[
E[Z^2]=\int_{-\infty}^{\infty} z^2 f_Z(z)\,dz
= \frac{1}{7}\int_{1}^{2} z^2\,dz + \frac{3}{7}\int_{5}^{7} z^2\,dz
= \frac{1}{7}\cdot\frac{2^3-1^3}{3} + \frac{3}{7}\cdot\frac{7^3-5^3}{3}
= \frac{661}{21}.
\]

\[
\mathrm{Var}(Z)=E[Z^2]-(E[Z])^2=\frac{661}{21}-\left(\frac{75}{14}\right)^2=\frac{1633}{588}.
\]

\finalanswer{E[Z]=\frac{75}{14},\ \mathrm{Var}(Z)=\frac{1633}{588}}
\end{solution}

\newpage

\problem{3.18}
Let $X$ be a normal random variable with mean $3$ and variance $4$.
\begin{enumerate}[label=(\alph*)]
\item Find the probability $P(2 < X < 6)$.
\item Find the value $c$ such that $P(X > c) = 0.33$.
\item Find $E(X^2)$.
\end{enumerate}
Hint. You can integrate with the density function, but it is quicker to relate $E(X^2)$ to the mean and variance.

\begin{solution}
\begin{scratch}
$\mu=3$, $\sigma=2$. Standardize: $Z=\frac{X-\mu}{\sigma}\sim N(0,1)$ with CDF $\Phi$.
(a) Compute $\Phi(1.5)-\Phi(-0.5)$.
(b) Solve $\Phi\!\left(\frac{c-3}{2}\right)=0.67$, so $c=3+2z_{0.67}$.
(c) Use $E(X^2)=\mathrm{Var}(X)+(E[X])^2$.
\end{scratch}

\textbf{(a)}
\[
P(2<X<6)=P\!\left(\frac{2-3}{2}<Z<\frac{6-3}{2}\right)=P\left(-\tfrac{1}{2}<Z<\tfrac{3}{2}\right)
=\Phi(1.5)-\Phi(-0.5).
\]
Using $\Phi(1.5)\approx0.9332$ and $\Phi(0.5)\approx0.6915$,
\[
P(2<X<6)\approx 0.9332-(1-0.6915)=0.6247.
\]

\textbf{(b)} We need $c$ such that
\[
0.33=P(X>c)=P\!\left(Z>\frac{c-3}{2}\right)=1-\Phi\!\left(\frac{c-3}{2}\right).
\]
Thus $\Phi\!\left(\frac{c-3}{2}\right)=0.67$. Let $z_{0.67}$ satisfy $\Phi(z_{0.67})=0.67$; from the table $z_{0.67}\approx0.44$.
Hence $c=3+2z_{0.67}\approx 3+2(0.44)=3.88$.

\textbf{(c)}
\[
E(X^2)=\mathrm{Var}(X)+(E[X])^2=4+3^2=13.
\]

\finalanswer{P(2<X<6)\approx0.6247,\ c\approx3.88,\ E(X^2)=13}
\end{solution}

\newpage

\problem{3.23}
Ten thousand people each buy a lottery ticket. Each lottery ticket costs \$1. 100 people are chosen as winners. Of those 100 people, 80 will win \$2, 19 will win \$100, and one lucky winner will win \$7000. Let $X$ denote the profit ($\text{profit} = \text{winnings} - \text{cost}$) of a randomly chosen player of this game.
\begin{enumerate}[label=(\alph*)]
\item Give both the possible values and probability mass function for $X$.
\item Find $P(X \ge 100)$.
\item Compute $E[X]$ and $\mathrm{Var}(X)$.
\end{enumerate}

\begin{solution}
\begin{scratch}
Possible winnings: $0,2,100,7000$. Profit $X=\text{winnings}-1$ so values are $-1,1,99,6999$.
Probabilities: 9900 nonwinners, 80 win $\$2$, 19 win $\$100$, 1 win $\$7000$ out of 10,000.
Compute $E[X]$, $E[X^2]$, then $\mathrm{Var}(X)=E[X^2]-(E[X])^2$.
\end{scratch}

\textbf{(a)} There are 10{,}000 tickets and 100 winners. Hence
\[
P(X=-1)=\frac{9900}{10000}=\frac{99}{100},\quad
P(X=1)=\frac{80}{10000}=\frac{1}{125},
\]
\[
P(X=99)=\frac{19}{10000},\quad
P(X=6999)=\frac{1}{10000}.
\]

\textbf{(b)} Since $99<100$, the only value with $X\ge100$ is $6999$:
\[
P(X\ge100)=P(X=6999)=\frac{1}{10000}.
\]

\textbf{(c)}
\[
E[X]=(-1)\frac{99}{100}+1\cdot\frac{1}{125}+99\cdot\frac{19}{10000}+6999\cdot\frac{1}{10000}
= -\frac{47}{500}\approx-0.094.
\]
\[
E[X^2]=1\cdot\frac{99}{100}+1\cdot\frac{1}{125}+99^2\cdot\frac{19}{10000}+6999^2\cdot\frac{1}{10000}
=\frac{245911}{50}=4918.22.
\]
\[
\mathrm{Var}(X)=E[X^2]-(E[X])^2=\frac{245911}{50}-\left(\frac{47}{500}\right)^2\approx 4918.211.
\]

\finalanswer{E[X]=-\frac{47}{500},\ \mathrm{Var}(X)\approx 4918.211,\ P(X\ge100)=\frac{1}{10000}}
\end{solution}

\newpage

\problem{3.28}
There are five closed boxes on a table. Three of the boxes have nice prizes inside. The other two do not. You open boxes one at a time until you find a prize. Let $X$ be the number of boxes you open.
\begin{enumerate}[label=(\alph*)]
\item Find the probability mass function of $X$.
\item Find $E[X]$.
\item Find $\mathrm{Var}(X)$.
\item Suppose the prize inside each of the three good boxes is \$100, but each empty box you open costs you \$100. What is your expected gain or loss in this game?
\end{enumerate}
Hint. Express the gain or loss as a function of $X$.

\begin{solution}
\begin{scratch}
$X$ = number of boxes opened until first prize. Possible values: $1,2,3$.
P($X=1$): first box good. P($X=2$): first empty, then good. P($X=3$): two empties, then good.
Then compute $E[X]$, $E[X^2]$, and $\mathrm{Var}(X)=E[X^2]-(E[X])^2$.
For (d) let gain/loss $W=100(2-X)$ and use $E[aX+b]=aE[X]+b$.
\end{scratch}

\textbf{(a)}
\[
P(X=1)=\frac{3}{5},\quad
P(X=2)=\frac{2}{5}\cdot\frac{3}{4}=\frac{3}{10},\quad
P(X=3)=\frac{2}{5}\cdot\frac{1}{4}\cdot\frac{3}{3}=\frac{1}{10}.
\]

\textbf{(b)}
\[
E[X]=1\cdot\frac{3}{5}+2\cdot\frac{3}{10}+3\cdot\frac{1}{10}=\frac{3}{2}.
\]

\textbf{(c)}
\[
E[X^2]=1^2\cdot\frac{3}{5}+2^2\cdot\frac{3}{10}+3^2\cdot\frac{1}{10}=\frac{27}{10},
\]
\[
\mathrm{Var}(X)=E[X^2]-(E[X])^2=\frac{27}{10}-\left(\frac{3}{2}\right)^2=\frac{9}{20}.
\]

\textbf{(d)} Each empty box costs \$100 and the prize is \$100, so the gain/loss is
\[
W=100-100(X-1)=100(2-X).
\]
By linearity,
\[
E[W]=100(2-E[X])=100\left(2-\frac{3}{2}\right)=50.
\]

\finalanswer{P(X=1,2,3)=\left(\frac{3}{5},\frac{3}{10},\frac{1}{10}\right),\ E[X]=\frac{3}{2},\ \mathrm{Var}(X)=\frac{9}{20},\ E[W]=50}
\end{solution}

\newpage

\problem{3.66}
Let $X$ be a normal random variable with mean $8$ and variance $3$. Find the value of $\alpha$ such that $P(X > \alpha) = 0.15$.

\begin{solution}
\begin{scratch}
$\mu=8$, $\sigma=\sqrt{3}$. Let $Z=\frac{X-8}{\sqrt{3}}\sim N(0,1)$.
Solve $P(Z>(\alpha-8)/\sqrt{3})=0.15$, so $\Phi((\alpha-8)/\sqrt{3})=0.85$.
\end{scratch}

\[
0.15=P(X>\alpha)=P\!\left(Z>\frac{\alpha-8}{\sqrt{3}}\right)=1-\Phi\!\left(\frac{\alpha-8}{\sqrt{3}}\right).
\]
Thus $\Phi\!\left(\frac{\alpha-8}{\sqrt{3}}\right)=0.85$. From the table, $z_{0.85}\approx1.04$.
Therefore
\[
\alpha=8+\sqrt{3}\,z_{0.85}\approx 8+\sqrt{3}(1.04)\approx 9.80.
\]

\finalanswer{\alpha=8+\sqrt{3}\,z_{0.85}\approx 9.80}
\end{solution}

\newpage

\problem{3.71}
My bus is scheduled to depart at noon. However, in reality the departure time varies randomly, with average departure time 12 o'clock noon and a standard deviation of 6 minutes. Assume the departure time is normally distributed. If I get to the bus stop 5 minutes past noon, what is the probability that the bus has not yet departed?

\begin{solution}
\begin{scratch}
Let $X$ be the departure time in minutes relative to noon. Then $X\sim N(0,36)$ and $Z=X/6\sim N(0,1)$.
We need $P(X>5)=P(Z>5/6)$ and relate to $\Phi$.
\end{scratch}

\[
P(\text{bus not yet departed})=P(X>5)=P\!\left(Z>\frac{5}{6}\right)=1-\Phi\!\left(\frac{5}{6}\right).
\]
Using $\Phi(0.83)\approx0.7967$,
\[
P(X>5)\approx 1-0.7967=0.2033.
\]

\finalanswer{P(X>5)=1-\Phi(5/6)\approx 0.2033}
\end{solution}


\end{document}
